\begingroup
\renewcommand{\arraystretch}{1.45}
\fancyhead[L]{\small BESTIMMUNG DER ANZIEHUNG}
\Needspace{32\baselineskip}
\AttractionArticle{3}
Es kommt nun darauf an, durch Einführung einer anderen Variablen anstelle von $E$ den Wurzelausdruck in eine einfachere Form zu bringen. Zu diesem Zweck werden wir setzen
\[
\begin{aligned}
&\cos E=\frac{\alpha+\alpha'\cos T+\alpha''\sin T}{\gamma+\gamma'\cos T+\gamma''\sin T},\\[6pt]
&
\sin E=\frac{\beta+\beta'\cos T+\beta''\sin T}{\gamma+\gamma'\cos T+\gamma''\sin T}
\end{aligned}
\]
Dabei sind jedoch die neun Koeffizienten $\alpha,\alpha',\alpha''$ usw. offenbar nicht völlig beliebig, sondern müssen gewissen Bedingungen genügen, die vor allem zu untersuchen sind. Zunächst bemerken wir, dass die Einsetzung dieselbe bleibt, wenn alle Koeffizienten mit demselben Faktor multipliziert werden, sodass man ohne Verlust an Allgemeinheit einem von ihnen einen festen Wert zuweisen, z. B. $\gamma=1$ setzen dürfte: Dennoch sollen der Eleganz halber alle neun noch eine Weile unbestimmt bleiben. Ferner weisen wir darauf hin, dass Werte ausge\SourceLeaf{344}%
schlossen werden müssen, bei denen $\alpha,\alpha',\alpha''$ oder $\beta,\beta',\beta''$ den $\gamma,\gamma'.\gamma''$ jeweils proportional wären: Andernfalls bliebe $E$ nicht mehr unbestimmt. Somit können $\gamma'\alpha''-\gamma''\alpha',\gamma''\alpha-\gamma\alpha'',\gamma\alpha'-\gamma'\alpha$ nicht zugleich verschwinden.

\Needspace{20\baselineskip}
Offenbar müssen die Koeffizienten $\alpha,\alpha',\alpha''$ usw. so beschaffen sein, dass identisch gilt
\[
\left.\begin{aligned}
 &(\alpha+\alpha'\cos T+\alpha''\sin T)^2\\[5pt]
+&(\beta+\beta'\cos T+\beta''\sin T)^2\\[5pt]
-&(\gamma+\gamma'\cos T+\gamma''\sin T)^2
\end{aligned}\right\}=0
\]
weshalb diese Funktion notwendig die Form
\[
\begin{aligned}
&k(\cos T^2+\sin T^2-1)
\end{aligned}
\]
haben muss. Hieraus erhalten wir sechs Bedingungsgleichungen
\[
\left.\begin{aligned}
-\alpha\alpha-\beta\beta+\gamma\gamma&=k\\[5pt]
-\alpha'\alpha'-\beta'\beta'+\gamma'\gamma'&=-k\\[5pt]
-\alpha''\alpha''-\beta''\beta''+\gamma''\gamma''&=-k\\[5pt]
-\alpha'\alpha''-\beta'\beta''+\gamma'\gamma''&=0\\[5pt]
-\alpha''\alpha-\beta''\beta+\gamma''\gamma&=0\\[5pt]
-\alpha\alpha'-\beta\beta'+\gamma\gamma'&=0
\end{aligned}\right\}\quad\mathrm{(I)}
\]
\Needspace{27\baselineskip}
Von diesen Gleichungen hängen mehrere andere ab, deren Entwicklung sich lohnen wird. Setzt man der Kürze halber
\[
\begin{aligned}
&\alpha\beta'\gamma''+\alpha'\beta''\gamma+\alpha''\beta\gamma'
-\alpha\beta''\gamma'-\alpha'\beta\gamma''-\alpha''\beta'\gamma
\\[6pt]
&\qquad{}=\varepsilon\quad\mathinner{.\,.\,.\,.\,.\,.\,.\,.}\quad\text{(II)}
\end{aligned}
\]
so werden durch Verbindung der Gleichungen (I) leicht die folgenden neun hergeleitet:
\[
\left.\begin{aligned}
\varepsilon\alpha&=-k(\beta'\gamma''-\gamma'\beta'')\\[5pt]
\varepsilon\beta&=-k(\gamma'\alpha''-\alpha'\gamma'')\\[5pt]
\varepsilon\gamma&=+k(\alpha'\beta''-\beta'\alpha'')\\[5pt]
\varepsilon\alpha'&=+k(\beta''\gamma-\gamma''\beta)\\[5pt]
\varepsilon\beta'&=+k(\gamma''\alpha-\alpha''\gamma)\\[5pt]
\varepsilon\gamma'&=-k(\alpha''\beta-\beta''\alpha)\\[5pt]
\varepsilon\alpha''&=+k(\beta\gamma'-\gamma\beta')\\[5pt]
\varepsilon\beta''&=+k(\gamma\alpha'-\alpha\gamma')\\[5pt]
\varepsilon\gamma''&=-k(\alpha\beta'-\beta\alpha')
\end{aligned}\right\}\quad\mathrm{(III)}
\]
\Needspace{38\baselineskip}
\SourceLeaf{345}
Aus den ersten drei dieser Gleichungen leiten wir wiederum die folgende her:
\[
\begin{aligned}
&\varepsilon\alpha(\beta'\gamma''-\gamma'\beta'')+\varepsilon\beta(\gamma'\alpha''-\alpha'\gamma'')+\varepsilon\gamma(\alpha'\beta''-\beta'\alpha'')\\[6pt]
&=-k(\beta'\gamma''-\gamma'\beta'')^2-k(\gamma'\alpha''-\alpha'\gamma'')^2\\[6pt]
&\qquad{}+k(\alpha'\beta''-\beta'\alpha'')^2
\end{aligned}
\]
der die folgende gleichwertig ist:
\[
\begin{aligned}
&\varepsilon\varepsilon=k(-\alpha'\alpha'-\beta'\beta'+\gamma'\gamma')(-\alpha''\alpha''-\beta''\beta''+\gamma''\gamma'')\\[6pt]
&\qquad{}-k(-\alpha'\alpha''-\beta'\beta''+\gamma'\gamma'')^2
\end{aligned}
\]
Diese geht mithilfe der Gleichungen 2, 3, 4 in (I) in die folgende über:
\[
\begin{aligned}
&\varepsilon\varepsilon=k^3\;\ldotp\;\ldotp\;\ldotp\;\ldotp\;\ldotp\;\ldotp\;\ldotp\;\ldotp\;\ldotp\;\ldotp\;\ldotp\;\text{(IV)}
\end{aligned}
\]
Ebenso leicht werden aus den Gleichungen (I) die folgenden hergeleitet:
\[
\left.\begin{aligned}
(\beta'\gamma''-\gamma'\beta'')^2&=-k(k-\alpha'\alpha'-\alpha''\alpha'')\\[5pt]
(\gamma'\alpha''-\alpha'\gamma'')^2&=-k(k-\beta'\beta'-\beta''\beta'')\\[5pt]
(\alpha'\beta''-\beta'\alpha'')^2&=+k(k+\gamma'\gamma'+\gamma''\gamma'')\\[5pt]
(\beta''\gamma-\gamma''\beta)^2&=+k(k+\alpha\alpha-\alpha''\alpha'')\\[5pt]
(\gamma''\alpha-\alpha''\gamma)^2&=+k(k+\beta\beta-\beta''\beta'')\\[5pt]
(\alpha''\beta-\beta''\alpha)^2&=-k(k-\gamma\gamma+\gamma''\gamma'')\\[5pt]
(\beta\gamma'-\gamma\beta')^2&=+k(k+\alpha\alpha-\alpha'\alpha')\\[5pt]
(\gamma\alpha'-\alpha\gamma')^2&=+k(k+\beta\beta-\beta'\beta')\\[5pt]
(\alpha\beta'-\beta\alpha')^2&=-k(k-\gamma\gamma+\gamma'\gamma')
\end{aligned}\right\}\quad\text{(V)}
\]

\Needspace{14\baselineskip}
Als Beispiel schreiben wir die Entwicklung der ersten hin, nach deren Muster die übrigen leicht zu bilden sein werden. Die Gleichungen 4, 2, 3 in (I) liefern nämlich
\[
\begin{aligned}
&(\gamma'\gamma''-\beta'\beta'')^2-(\gamma'\gamma'-\beta'\beta')(\gamma''\gamma''-\beta''\beta'')\\[6pt]
&\qquad{}=\alpha'\alpha'\alpha''\alpha''-(\alpha'\alpha'-k)(\alpha''\alpha''-k)
\end{aligned}
\]
Diese Gleichung ergibt entwickelt unmittelbar die erste in (V) selbst.

\Needspace{26\baselineskip}
Aus diesen Gleichungen (V) schließen wir, dass der Wert $k=0$ in unserer Untersuchung nicht zulässig ist; denn dann würden alle neun Größen $\beta'\gamma''-\nobreak{}\gamma'\beta''$ usw. notwendig verschwinden, d. h. die Koeffizienten $\alpha$, $\alpha'$, $\alpha''$ würden sowohl den $\beta$, $\beta'$, $\beta''$ als auch den $\gamma$, $\gamma'$, $\gamma''$ proportional werden. Daher kann auch wegen der Gleichung IV die Größe $\varepsilon$ nicht verschwinden; folglich muss $k$ notwendig eine positive Größe sein, sofern alle Koeffizienten $\alpha$, $\alpha'$, $\alpha''$ usw. reell sein sollen. Verbindet man die ersten drei Gleichungen in (III) mit den ersten drei in (V), so ergeben sich die folgenden neuen, die offenbar davon abhängen, dass der Wert von $k$ nicht verschwindet:
\SourceLeaf{346}
\[
\left.\begin{aligned}
\alpha\alpha-\alpha'\alpha'-\alpha''\alpha''&=-k\\[5pt]
\beta\beta-\beta'\beta'-\beta''\beta''&=-k\\[5pt]
\gamma\gamma-\gamma'\gamma'-\gamma''\gamma''&=+k
\end{aligned}\right\}\quad\text{(VI)}
\]
\Needspace{20\baselineskip}
Die Verbindung der übrigen würde dieselben ergeben. Diesen fügen wir schließlich die folgenden drei hinzu:
\[
\left.\begin{aligned}
\beta\gamma-\beta'\gamma'-\beta''\gamma''&=0\\[5pt]
\gamma\alpha-\gamma'\alpha'-\gamma''\alpha''&=0\\[5pt]
\alpha\beta-\alpha'\beta'-\alpha''\beta''&=0
\end{aligned}\right\}\quad\text{(VII)}
\]
Sie werden leicht aus den Gleichungen III hergeleitet; z. B. liefern die zweite, fünfte und achte:
\[
\begin{aligned}
&\varepsilon\beta\gamma-\varepsilon\beta'\gamma'-\varepsilon\beta''\gamma''=-k\gamma(\gamma'\alpha''-\alpha'\gamma'')\\[6pt]
&\qquad{}-k\gamma'(\gamma''\alpha-\alpha''\gamma)-k\gamma''(\gamma\alpha'-\alpha\gamma')=0
\end{aligned}
\]
Offenbar hängen auch diese Gleichungen vom Ausschluss des Wertes $k=0$ ab\HistoricalNote{25}.

\Needspace{36\baselineskip}
Da, wie wir schon oben bemerkt haben, alle Koeffizienten $\alpha$, $\alpha'$, $\alpha''$ usw. mit demselben Faktor multipliziert werden dürfen, wobei der Wert von $k$ mit dem Quadrat desselben Faktors multipliziert hervorgehen wird, werden wir von nun an stets annehmen
\[
\begin{aligned}
&k=1
\end{aligned}
\]
Dabei wird notwendig auch entweder $\varepsilon=+1$ oder $\varepsilon=-1$ sein. Offenbar müssen also die neun Koeffizienten $\alpha$, $\alpha'$, $\alpha''$ usw., zwischen denen sechs Bedingungsgleichungen bestehen, auf drei voneinander unabhängige Größen zurückführbar sein, was sich am bequemsten durch drei Winkel auf folgende Weise ausführen lässt:
\[
\begin{aligned}
&\alpha=\cos L\operatorname{tang}N\\[6pt]
&\beta=\sin L\operatorname{tang}N\\[6pt]
&\gamma=\sec N\\[6pt]
&\alpha'=\cos L\cos M\sec N\pm\sin L\sin M\\[6pt]
&\beta'=\sin L\cos M\sec N\mp\cos L\sin M\\[6pt]
&\gamma'=\cos M\operatorname{tang}N\\[6pt]
&\alpha''=\cos L\sin M\sec N\mp\sin L\cos M\\[6pt]
&\beta''=\sin L\sin M\sec N\pm\cos L\cos M\\[6pt]
&\gamma''=\sin M\operatorname{tang}N
\end{aligned}
\]

\SourceLeaf{347}\noindent Dabei beziehen sich von den doppelten Vorzeichen die oberen auf den Fall $\varepsilon=+1$, die unteren auf den Fall $\varepsilon=-1$. Die analytische Behandlung lässt sich jedoch zum größten Teil ohne Verwendung dieser Winkel eleganter durchführen. Im Übrigen wäre es nicht schwierig, sowohl diesen Winkeln als auch den übrigen in dieser Untersuchung vorkommenden Hilfsgrößen eine geometrische Bedeutung zuzuweisen; diese für unseren Zweck nicht notwendige Deutung überlassen wir jedoch dem kundigen Leser zur Erläuterung.


\clearpage
\endgroup
